% ===================== 练习：二十题 =====================
\clearpage
\phantomsection\label{sec:practice}
\hdr{练习}{二十题：课件作业与自编题}

正文各节讲完一个方法，这里集中练。二十题分五组，来源不同，先看清楚再动笔：

\begin{itemize}
  \item \textbf{P1--P4}（§7.1）：2026-11-02 旧课件页面上留下的四个问题，其中 P4 是按课件 p14 公式编的数字练习。这四题当时随旧稿给过详解，属于复习。
  \item \textbf{P5--P8}（§7.4）：课件 Chapter 7 p19 的四道定解条件作业。2026-11-09 已出过中英详解，同样是复习；但符号仍是错得最多的地方，值得重做。
  \item \textbf{P9--P12}（§8.1）：课件 Chapter 8 p13 的四道分离变量作业。2026-11-13 给过分步详解。
  \item \textbf{P13--P16}（§8.2）：自编四题，覆盖圆内、环域、半圆盘与\"哪些函数能作解\"的判断。
  \item \textbf{P17--P20}（§8.3）：P17、P18 是课件 p6 的作业 5、7；P19、P20 自编，练 Fourier--Bessel 展开与 $\mu$ 的判据。
\end{itemize}

每题只给题干与考点，步骤全在附录（P1--P20 同号）。建议合上讲义先写一遍，再翻附录对照；
尤其 P7、P8 的符号与 P10 的 $\omega_n$，是往年错得最多的两处。

\hdrsub{A}{§7.1 的四道思考题}

\practice{P1}{课件 p8 思考题}{同一种弦，\textbf{把弦绷紧}与\textbf{换一根更粗的弦}，音高分别怎么变？请用波速公式解释。}{$a=\sqrt{T/\rho}$ 中 $T$ 与 $\rho$ 各自的作用方向；音高由频率决定}{ans:P1}

\practice{P2}{课件 p13 提问}{杆的纵振动加上「单位长度外力」$F(x,t)$ 后，方程变成什么样？若外力是杆的自重（密度 $\rho$、重力加速度 $g$、方向沿 $-x$），方程又怎么写？}{把外力并进非齐次项时要除以单位长度的质量；恒力的非齐次项是常数}{ans:P2}

\practice{P3}{课件 p3 提问}{独立推导一维热传导方程，并说明 $a^2$ 与 $f$ 各由哪些物理量组成。条件：截面积 $S$、体密度 $\rho$、比热容 $c$、热导率 $\kappa$、单位体积生热率 $F$。}{能量守恒加傅里叶定律；$\rho c$ 与 $\kappa$ 的出现位置}{ans:P3}

\practice{P4}{按课件 p14 公式编排}{膜的面密度 $\rho=0.020\ \unit{kg\cdot m^{-2}}$，张力 $T=2.0\ \unit{N\cdot m^{-1}}$。(1) 求波速 $a$；(2) 若膜上还受均匀的单位面积外力 $f=0.50\ \unit{m\cdot s^{-2}}$，写出方程；(3) 这个 $f$ 的量纲为什么是加速度？}{$a=\sqrt{T/\rho}$ 的量纲核对；非齐次项是加速度而不是力}{ans:P4}

\hdrsub{B}{§7.4 的四道定解条件题（课件 Ch7 p19）}

\practice{P5}{课件 Ch7 p19 第 1 题}{长 $\ell$ 的杆，左端固定；右端受一个沿外法向、大小为 $F(t)=F_0\sin\omega t$ 的力。写出这个问题的边界条件。}{第一类与第二类边界条件的区分：题目给的是值还是力}{ans:P5}

\practice{P6}{课件 Ch7 p19 第 2 题}{长 $\ell$ 的弦两端固定。$t=0$ 时在 $x=h$ 处受一个外力 $F_0$。写出这个问题的初始条件。}{时间二阶问题要几个初始条件；集中冲量与集中位移在数学表示上的区别}{ans:P6}

\practice{P7}{课件 Ch7 p19 第 3 题}{长 $\ell$ 的杆，两端各受一个大小 $F_0$ 的法向力，方向都指向杆内。写出边界条件。}{两端受压时第二类条件的符号，以及两个端点的方向是否一致}{ans:P7}

\practice{P8}{课件 Ch7 p19 第 4 题}{长 $\ell$ 的杆，两端各有热流 $q_0$ 流入杆内。写出边界条件。}{热流边界的方向判定：流入、流出与外法向的关系}{ans:P8}

\hdrsub{C}{§8.1 的四道分离变量题（课件 Ch8 p13）}

\practice{P9}{课件 Ch8 p13 第 1 题}{长 $\ell$ 的弦两端固定，初始位移为 $\varphi(x)=hx\ (0<x<\ell/2)$、$\varphi(x)=h(\ell-x)\ (\ell/2<x<\ell)$，初始速度为零，$h$ 为常数。求 $u(x,t)$。}{两端固定弦的特征值与 $\sin$ 级数；三角初位移的系数积分}{ans:P9}

\practice{P10}{课件 Ch8 p13 第 2 题}{长 $\ell$ 的杆左端固定、右端自由；初始位移为零，初始速度 $\psi(x)=bx/\ell$。求 $u(x,t)$。}{一端固定一端自由的特征值问题；初速度给定时的系数定法}{ans:P10}

\practice{P11}{课件 Ch8 p13 第 3 题}{长 $\ell$ 的杆两端温度为零，初始温度 $u(x,0)=bx(\ell-x)/\ell^2$。求温度分布 $u(x,t)$。}{输运方程的分离变量；初值的对称性如何减少非零项}{ans:P11}

\practice{P12}{课件 Ch8 p13 第 4 题}{矩形 $0<x<a$、$0<y<b$，求 Laplace 方程的解，边界条件为 $u(0,y)=u(a,y)=0$、$u(x,0)=A(1-x/a)$、$u(x,b)=0$。}{二维稳态的两次分离；非齐次边界位置如何决定双曲函数的形式}{ans:P12}

\hdrsub{D}{§8.2 的四道极坐标题（自编）}

\practice{P13}{自编，仿课件 p3 的圆内类型}{单位圆盘 $r<1$ 内解 $\nabla^2u=0$，圆周上 $u(1,\varphi)=\sin^2\varphi$。(1) 求 $u(r,\varphi)$；(2) 把它写成直角坐标 $x,y$ 的形式，并说明它为什么满足 $\nabla^2u=0$。}{边界值的余弦级数展开；圆内自然边界条件筛项}{ans:P13}

\practice{P14}{自编，仿课件 p5 的环域类型}{同轴圆筒之间 $a<r<b$：内筒 $u(a,\varphi)=0$，外筒 $u(b,\varphi)=U_0$（常数）。(1) 求 $u$；(2) 求径向场 $E_r$ 与内筒表面电荷密度 $\sigma_a$；(3) 求单位长度电容，并与平行板电容的极限对照。}{$m=0$ 分支；$\ln r$；两个 $r$ 条件定两个常数}{ans:P14}

\practice{P15}{自编，仿课件 p5 的类型但角向条件改为固定端点}{上半圆盘 $r<a,\ 0<\varphi<\pi$ 内解 $\nabla^2u=0$，边界条件为 $u(r,0)=u(r,\pi)=0$（$r<a$）与 $u(a,\varphi)=U_0$（$0<\varphi<\pi$）。(1) 写出 $\varphi$ 方向的特征值与特征函数；(2) 求 $u$ 的级数解，写出前三项；(3) 说明级数在 $r=a$ 上收敛到什么，以及两个角点 $(a,0)$、$(a,\pi)$ 处发生了什么。}{角向为固定端点条件时的谱与特征函数；奇偶项取舍}{ans:P15}

\practice{P16}{自编，判断类型}{下列四个函数都在 $r>0$ 上满足二维 Laplace 方程（可自行验算）。问：哪几个可以作为单位圆盘 $r<1$ 内 Laplace 方程的解？对每一个说明是与不是的理由，理由要落到具体条件上。\ (a) $r^{-1}\cos\varphi$；\ (b) $\ln r$；\ (c) $r^2\sin2\varphi$；\ (d) $r^{1/2}\cos(\varphi/2)$。}{通解的两道筛子：$\varphi$ 方向的单值与 $r$ 方向端点的有界性}{ans:P16}

\hdrsub{E}{§8.3 的四道柱坐标与贝塞尔题}

\practice{P17}{课件 p6 作业 5}{在半径 $\rho_0$ 的圆内解 $\nabla^2u=0$，边界条件分别为：(1) $u\big|_{r=\rho_0}=A\cos\varphi$；(2) $u\big|_{r=\rho_0}=A+B\sin\varphi$，其中 $A,B$ 为常数。分别求 $u$。}{圆内 Laplace 的 $r^m$ 项；边界值就是一个谐波时系数直接读出}{ans:P17}

\practice{P18}{课件 p6 作业 7}{内半径 $\rho_1$、外半径 $\rho_2$ 的环域内解稳态热传导 $\nabla^2u=0$；内圆温度 $u\big|_{r=\rho_1}=u_1\cos2\varphi$，外圆温度 $u\big|_{r=\rho_2}=u_2\sin\varphi$。求温度分布。}{环域保留 $r^{\pm m}$；两个不同阶的谐波各自独立定系数}{ans:P18}

\practice{P19}{自编，Fourier--Bessel 展开}{在 $0<r<a$ 上用 $J_0(x_n^{(0)}r/a)$ 展开 $f(r)=1-r^2/a^2$。(1) 写出系数公式；(2) 算出前三个系数（保留三位有效数字）；(3) 说明为什么这个级数在 $r=a$ 处给出 $0$，与 $f(a)=0$ 一致。}{Fourier--Bessel 系数公式；$\int_0^z tJ_0(t)\dd t=zJ_1(z)$}{ans:P19}

\practice{P20}{自编，$\mu$ 的判据}{下列四个柱坐标问题，分别应取 $\mu>0$、$\mu<0$ 还是 $\mu=0$？说明判据落在哪一组边界上。(a) 有限高圆柱，侧面 $u(a,\varphi,z)=0$，上下底给定非零温度；(b) 有限高圆柱，上下底 $u=0$，侧面给定温度；(c) 无限长圆柱（无 $z$ 边界），侧面 $u(a,\varphi)=f(\varphi)$；(d) 有限高圆柱，侧面绝热 $\partial_ru(a,\varphi,z)=0$，上下底给定温度，且温度与 $\varphi$ 无关。}{「哪组边界齐次」判据；$\mu=0$ 的轴对称欧拉分支}{ans:P20}

